\section*{Введение}
\addcontentsline{toc}{section}{Введение}
\label{sec:introduction}
\draft{Написать после основной части: задача турнира, особенности меняющейся геометрии, температурной истории и калибровки; затем кратко перечислить полученные результаты.}

Температурная история осадочного бассейна определяет условия преобразования органического вещества в породах. Одним из её показателей является отражательная способность витринита, характеризующая его зрелость. В отличие от современной температуры, описывающей текущее тепловое состояние, отражательная способность зависит от воздействия температуры за время существования материала. Поэтому совместное использование температурных измерений и данных о зрелости позволяет учитывать разные стороны истории прогрева разреза.

В постановке задачи Турнира~\cite{problem2026} геометрическая история осадочного бассейна задаётся последовательностью конфигураций, отражающих осадконакопление, движение слоёв и разломные смещения. Решение температурной задачи требует учитывать изменение геометрии, неоднородность литотипов, пористость и зависимость теплофизических свойств от температуры. При вычислении отражательной способности дополнительно необходимо сохранять температурную историю каждого материала и учитывать момент его отложения. В обратной задаче такой прямой расчёт повторяется для разных граничных условий, что предъявляет требования одновременно к численной точности и вычислительным затратам.

Цель работы~--- разработать и проверить численный алгоритм определения температурного поля и отражательной способности витринита, а также восстановления базального теплового потока по современным скважинным наблюдениям. Для этого последовательно рассматриваются математическая модель среды, численные методы решения прямой и обратной задач и их программная реализация. Достоверность результатов оценивается на задачах с известным решением и при сопоставлении с предоставленными данными. При этом точность расчёта полей, согласование наблюдений и точность восстановления потока рассматриваются как отдельные характеристики.

В работе используются зависимости теплофизических свойств, применяемые в бассейновом моделировании~\cite{hantschel2009}, и кинетическая модель EASY\%Ro~\cite{sweeney1990}. Для температурной задачи построена конечно-объёмная аппроксимация на подвижной колоночной сетке с многоточечным вычислением потоков. Обратная задача решается минимизацией суммы квадратов нормированных невязок температуры и отражательной способности при ограничениях на монотонность и диапазон базального потока. Проверка включает изготовленные решения, восстановление известного потока по синтетическим наблюдениям и сопоставление прямого расчёта с предоставленным эталонным решением Модели~1.

Отчёт состоит из пяти основных разделов. В разделе~\ref{sec:problem} приведены математическая модель среды и постановки прямой и обратной задач. Раздел~\ref{sec:algorithm} посвящён численным алгоритмам, а раздел~\ref{sec:implementation}~--- их программной реализации и вычислительным затратам. В разделе~\ref{sec:verification} исследуются численная точность и восстановление параметров на задачах с известным решением. В разделе~\ref{sec:results} рассмотрены расчёты на предоставленных данных, сравнение с эталонным решением и согласование измерений в скважинах.
